\section{Binary Front Ends and Injection}
\label{sec:frontends}

\subsection{ELF64 parsing and validation}

The ELF front end accepts 64-bit little-endian x86-64 executables and shared
objects. The parser validates the ELF header, class, data encoding, machine,
program-header table, section-header table, section bounds, load-segment flags,
and the presence of an executable range. Extended numbering (\code{PN_XNUM},
\code{SHN_XINDEX}) and unsupported object types fail closed. Function symbols
are read from \code{.symtab} or \code{.dynsym}; fully stripped inputs require
marker or explicit-range selection.

The editor never relocates the target. It appends new bytes and rewrites the
program-header table in a new read-only page. This avoids moving existing
segments and preserves the host loader's expectations.

\subsection{Selector resolution}

A protection run resolves one or more functions using one of four selectors:

\begin{itemize}
  \item an exact symbol name (\verb|--function|);
  \item a unique symbol substring (\verb|--function-contains|);
  \item an explicit executable RVA and size (\verb|--range|);
  \item raw byte markers (\verb|--marker-begin|/\verb|--marker-end|).
\end{itemize}

The resolved range must lie in an executable segment, must not overlap another
selected range or an already-patched entry, and must decode completely within the
supported subset. Marker mode is intentionally raw: the user is responsible for
placing the markers so that the enclosed bytes form valid machine code.

\subsection{ELF injection geometry}

The appended layout is:

\begin{center}
\small
\begin{tabular}{@{}lll@{}}
\toprule
Region & Mapping & Content \\
\midrule
Original bytes & original segments & unchanged host image \\
Program-header page & \code{R} & relocated \code{PT_PHDR} table \\
VM segment & \code{R+X} & engine, tables, payload, entry thunks \\
\code{.shstrtab} & section metadata & regenerated names \\
Section-header table & section metadata & regenerated entries \\
\bottomrule
\end{tabular}
\end{center}

One new read-only \code{PT_LOAD} maps the relocated program-header table, and one
new read-execute \code{PT_LOAD} maps the VM. The \code{PT_PHDR} entry is updated
to the new table. The section-header table and section-name string table are
regenerated; \code{.sirus} and \code{.sirus.payload} sections are created even
for inputs that lacked a name table.

The injected VM is linked at a chosen virtual address with a custom linker
script. The blob is a single read-execute segment with no writable load, no
relocations, and no undefined symbols. The validator rejects any blob that
violates these properties, any missing entry symbol, any entry outside the
segment, and any payload/seed/hash mismatch.

\subsection{Entry patching and body scrubbing}

The patcher writes \code{E9 <rel32>} at the target entry, or immediately after a
preserved \code{F3 0F 1E FA} (\code{endbr64}) sequence. If the displacement does
not fit in 32 bits, the build fails. The remainder of the original function body
is overwritten with deterministic random bytes derived from the build seed,
unless \verb|--no-scrub| is requested. The output is written to a temporary file
in the destination directory, permissions are copied from the input, and the file
is atomically renamed.

Preserving \code{endbr64} keeps the patched entry a valid indirect-branch target
under Intel CET. The jump itself preserves the host return address: the patched
function is entered by a normal call, and the VM thunk returns normally.

\subsection{PE32+ backend}

The Windows backend accepts PE32+ AMD64 images. It validates the DOS, COFF, and
optional headers, the machine type, section alignment, section bounds, RVA to file
mapping, checksum, and COFF symbol information. Injection appends one
read-execute \code{.sirus} section, updates the section count and
\code{SizeOfImage}, recomputes the PE checksum, patches one RVA-relative jump per
protected function, and overwrites the original body with keyed junk.

The Win64 runtime is cross-compiled with
\code{rustc --target x86_64-pc-windows-gnu}, linked with \code{lld-link}, and
embedded as one section. The generated shim includes a guard-page-safe stack
probe, preserves the Windows calling convention, and translates register and
stack arguments into the VM's internal array. Linked \code{.pdata} entries are
rebased into the section and merged into the target exception directory;
handler-bearing and chained unwind records fail closed.

The tested Windows fixtures cover a normal leaf, constant-time execution, a
hardened self-code checksum, a \code{/guard:cf} binary with indirect callers, a
DLL export loaded through \code{LoadLibrary}/\code{GetProcAddress}, and a TLS
callback plus \code{RCDATA} resource. The native Windows CI job, the Wine job,
and the MSVC \code{/guard:cf} job all pass on the release tree. Signing,
arbitrary imports, runtime exception propagation, and kernel anti-debug remain
out of scope.

\subsection{ABI shims and call semantics}

Table~\ref{tab:abi} summarizes the supported calling conventions.

\begin{table}[htbp]
\centering
\small
\caption{Supported integer ABI surface.}
\label{tab:abi}
\begin{tabularx}{\textwidth}{@{}l>{\raggedright\arraybackslash}X>{\raggedright\arraybackslash}X@{}}
\toprule
Property & System V AMD64 & Win64 \\
\midrule
Register arguments & 6: RDI, RSI, RDX, RCX, R8, R9 & 4: RCX, RDX, R8, R9 \\
Stack arguments & 6: 7th--12th & 6: 5th--10th \\
Integer result & RAX & RAX \\
Caller-saved GPRs & preserved by shim & preserved by shim \\
XMM preservation & XMM0--XMM15 saved by shim & XMM0--XMM5 volatile; XMM6--XMM15 callee-saved \\
Stack alignment & shim aligns before the VM call & guard-page-safe probe \\
Direct protected calls & bounded guest call stack & bounded guest call stack \\
External calls & rejected & rejected \\
\bottomrule
\end{tabularx}
\end{table}

The Linux shim saves the caller-saved GPRs and all sixteen XMM registers, records
\code{RBP}, aligns the stack defensively, copies stack arguments, and calls the
generated entry. The frame-pointer restore makes the shim robust when the host
caller has an unusual stack alignment, a case exposed by an eight-thread pthread
fixture. The Windows shim preserves the volatile GPRs and \code{XMM0}--\code{XMM5}
and relies on the callee-saved contract for \code{XMM6}--\code{XMM15}; the ABI
matrix exercises those registers, indirect callers, and a 32 KiB stack frame.

Direct calls between selected functions execute on a bounded guest call stack.
The VM validates the target entry state, pushes a return state, and returns to
the caller state on \code{RET}. Calls to an external or unresolved function fail
closed. The compiler can automatically lift resolvable direct callees with
\verb|--follow-calls|; unresolved targets still fail closed.

The VM uses only per-invocation stack state. It has no mutable global engine
state, so separate host threads do not share virtual registers or private
memory. An eight-thread regression and a SIGALRM handler reentrancy test exercise
this property. The payload, decode table, and seed are read-only after
initialization.

\subsection{Injection failure behavior}

Every front-end operation is transactional. If parsing, selection, lifting,
encoding, compilation, validation, or patching fails, the protector returns a
typed error and does not write a partial output. Existing output is replaced
atomically only after the new image is complete. This matters for build systems:
a failed protection run cannot leave a half-patched binary that appears valid.
